\documentclass[12pt]{article}
\usepackage{amsmath,amssymb,amsthm}
\usepackage[margin=1in]{geometry}

\newtheorem{theorem}{Theorem}
\newtheorem{lemma}[theorem]{Lemma}

\title{Problem 372: Pencils of Rays}
\author{Project Euler}
\date{}

\begin{document}
\maketitle

\section{Problem Statement}

Define $S(N) = \sum_{n=1}^{N} \lfloor(n+1)/2\rfloor\,\varphi(n)$, where $\varphi$ is Euler's totient function. Compute $S(10^{10})$.

\section{Mathematical Foundation}

\begin{theorem}[Totient Summation via M\"obius Inversion]
For any positive integer~$N$:
\[
\Phi(N) := \sum_{n=1}^{N} \varphi(n) = \frac{1}{2}\!\left(1 + \sum_{d=1}^{N} \mu(d)\left\lfloor \frac{N}{d}\right\rfloor\!\left(\left\lfloor \frac{N}{d}\right\rfloor + 1\right)\right).
\]
\end{theorem}

\begin{proof}
By M\"obius inversion, $\varphi(n) = \sum_{d \mid n} \mu(d)\,\frac{n}{d}$. Summing over~$n$:
\[
\Phi(N) = \sum_{d=1}^{N} \mu(d) \sum_{k=1}^{\lfloor N/d\rfloor} k
= \sum_{d=1}^{N} \mu(d)\,\frac{\lfloor N/d\rfloor\,(\lfloor N/d\rfloor+1)}{2}. \qedhere
\]
\end{proof}

\begin{lemma}[Weighted Totient Sum]
Define $F(N) = \sum_{n=1}^{N} f(n)\,\varphi(n)$ for a weight function~$f$. Using $\varphi = \mu * \mathrm{id}$:
\[
F(N) = \sum_{d=1}^{N} \mu(d) \sum_{k=1}^{\lfloor N/d\rfloor} f(dk)\cdot k.
\]
\end{lemma}

\begin{proof}
Substitute $\varphi(n) = \sum_{d \mid n}\mu(d)(n/d)$ and swap summation with $n = dk$.
\end{proof}

\begin{lemma}[Block Summation]
The values $\lfloor N/d \rfloor$ take at most $O(\sqrt{N})$ distinct values. For each block of consecutive $d$-values sharing the same quotient, the inner sum can be batched, reducing the total work to $O(\sqrt{N})$ for the tail.
\end{lemma}

\begin{proof}
For $q \le \sqrt{N}$, there is at most one block per $q$-value ($\le \sqrt{N}$ blocks). For $d \le \sqrt{N}$, there are at most $\sqrt{N}$ individual values. Hence $O(\sqrt{N})$ blocks total.
\end{proof}

The target sum $S(N) = \sum_{n=1}^{N}\lfloor(n+1)/2\rfloor\,\varphi(n)$ is computed by applying the weighted totient sum framework with $f(n) = \lfloor(n+1)/2\rfloor$, combined with sieve precomputation up to $\sim N^{2/3}$ and the hyperbola method for the tail.

\section{Algorithm}

\begin{verbatim}
function solve(N):
    B = ceil(N^(2/3))
    mu[1..B] = sieve_mobius(B)
    result = 0
    for d = 1 to B:
        if mu[d] == 0: continue
        Q = floor(N / d)
        result += mu[d] * inner_sum(d, Q)
    // Hyperbola tail for remaining blocks
    return result
\end{verbatim}

\section{Complexity Analysis}

\begin{itemize}
\item \textbf{Time:} $O(N^{2/3}\log\log N)$ for the sieve, plus $O(\sqrt{N})$ for the hyperbola tail. Dominant: $O(N^{2/3})$.
\item \textbf{Space:} $O(N^{2/3})$ for the sieve arrays.
\end{itemize}

\section{Answer}

\[
\boxed{301450082318807027}
\]

\end{document}
